% Additive teaching material. Deliberately use unnumbered environments so
% the retained research report keeps its theorem/equation/section numbering.
\begin{center}
\large A gradual course in thermal relaxation, from fractions to exact proof
\end{center}
\vspace{1em}
\noindent\textbf{Begin at the beginning.} The first part develops the needed
mathematics through 40 lessons, worked examples, and exercises with answers.
The second part retains the complete technical research report and its
reproducible certificates. Read them as a connected course, or use the
linked contents and reading map to move between intuition and proof.
\clearpage
\section*{How to read this unusually detailed report}
\addcontentsline{toc}{section}{How to read this unusually detailed report}
This report has two jobs. It teaches the mathematics needed to understand a
thermal relaxation counterexample, and it records the research arguments
precisely enough to check them. The first job requires patience. We begin with
fractions and probabilities, not with a matrix full of unfamiliar symbols.

The central question is simple to state. Prepare two copies of the same system,
one hotter than the other. Put both in contact with the same colder environment.
Can the hotter copy later be closer to equilibrium? In the models studied here,
the answer can be yes. A second question is more subtle: can every three-state
test miss a reversal in a four-state system? A certified example in this report
shows that this too can happen. Understanding why requires learning what the
state of the system is, what its dynamics do, what ``closer'' means, and how a
proof differs from a convincing numerical plot.

\paragraph{Starting knowledge.}
You need ordinary arithmetic and a willingness to read an equation one symbol
at a time. High-school algebra helps. Calculus and linear algebra are developed
as they become necessary. This is not a replacement for complete courses in
those subjects: it is a carefully chosen route through the parts this research
uses. No knowledge of Markov chains, spectral theory, statistical mechanics,
convex order, interval arithmetic, or proof assistants is assumed at the start.

\paragraph{The learning sequence.}
Lessons 1--9 build the language: logical claims, fractions, probabilities,
weighted averages, temperature, exponentials, Gibbs distributions, derivatives,
and covariance. Lessons 10--17 build the dynamics: rates, matrices,
equilibrium, eigenmodes, detailed balance, distances, and slow amplitudes.
Lessons 18--24 work through the counterexample and the structural explanation.
Lessons 25--30 bridge to classification, degeneracy, rigorous computation,
finite-time physics, and formal proof. Lessons 31--32 give exercises with
solutions and a map into the research text. Lessons 33--40 then slow down
the bridge to continuous-time certificates and physical models, deriving
random clocks, rigorous tails, target-time bounds and intermediate-state
dynamics through worked examples.

\paragraph{How much to read at once.}
For a first pass, read one lesson, reproduce its small numerical example, and
stop when the last paragraph makes sense. For a second pass, derive the displayed
identities without looking at the intermediate lines. For a research pass,
follow the references into Part II and inspect the associated certificates.
Being able to recite a conclusion is different from being able to reconstruct
its argument; the worked checks are intended to close that gap.

\paragraph{What the labels mean.}
An \emph{illustration} is chosen to make an idea easy to see. It need not exhibit
a Mpemba reversal. An \emph{identity} follows exactly from its stated assumptions.
A \emph{certified result} refers to the particular mathematical claim and
verification described in the research text. A \emph{diagnostic calculation}
provides evidence or intuition, but does not by itself certify a continuum of
parameters. The new teaching examples do not expand the scope of the existing
Lean formalization or turn an unfinished proof audit into a completed one.

\paragraph{Keep three questions separate throughout.}
What has been proved for an abstract finite-state model? What has been verified
for the particular parameter choices? What has been demonstrated in a physical
experiment? This report supplies mathematical and computational results for
models. The physical-model study explains its idealizations explicitly.

\clearpage
\tableofcontents
\clearpage
\part*{Part I. From basic arithmetic to the research frontier}
\addcontentsline{toc}{part}{Part I. From basic arithmetic to the research frontier}

\section*{Lesson 1. What would count as an answer?}
\addcontentsline{toc}{section}{1. What would count as an answer?}
\paragraph{Start with a familiar comparison.}
Suppose two cups of water start at different temperatures. Saying that one
``cools faster'' is incomplete. It might mean a steeper initial temperature
slope, a shorter time to a target temperature, or a smaller distance from the
final state after ten minutes. Those statements are different. A curve can
start steeply and later flatten. A system can cross one target first and a
different target second. We therefore choose a mathematical comparison before
trying to prove anything.

Here the system has finitely many states. Its condition is described by a list
of probabilities. Temperature is used to prepare that list; after the dynamics
begin, the list need not correspond to any single temperature. A bath is the
environment with fixed temperature. Both preparations evolve with the same
rules toward the same bath equilibrium. We compare distances between their
probability lists and that equilibrium.

\paragraph{An implication is a promise with a condition.}
Write $A\Rightarrow B$ for ``if $A$ holds, then $B$ holds.'' A proposed universal
test says that whenever certain three-state tests fail, a four-state reversal
cannot occur. To disprove such a statement, we must produce a system in which
the condition $A$ is true and the promised conclusion $B$ is false. Producing a
reversal alone is not enough: the system must also satisfy every condition of
the claim being tested.

There is an important asymmetry. One valid counterexample disproves a universal
claim. A million examples consistent with that claim do not prove it. Sampling
can miss a thin set, a boundary, or a region outside the sample box. Conversely,
failure to find a counterexample is not evidence that a particular proof is
correct. A proof must explain why every allowed case is covered.

\paragraph{Read the quantifiers.}
``There exists'' means that at least one object works. ``For every'' means that
none of the allowed objects may fail. In symbols,
\[
\exists M\;\text{with a reversal}
\qquad\hbox{and}\qquad
\forall M\;\text{no reversal}
\]
make very different claims. The letter $M$ will later denote a matrix of rates.
For now it is simply the name of the dynamical rule. A universal bound also
needs quantifiers: a bound for four states is not automatically a bound for all
numbers of states, and a bound for long-time amplitudes is not automatically a
bound for target times.

\paragraph{The proof obligations to remember.}
Our counterexample must have legitimate probabilities, distinct ordered
energies, a finite positive bath temperature, positive transition rates,
probability-conserving dynamics, detailed balance, the stated slow-mode
property, failure of all the prescribed triplet tests, and a strict reversal.
Each adjective has a job. Later lessons explain how each is checked.

\paragraph{Checkpoint.}
If a three-state test succeeds but a reversal is observed, have we disproved a
claim whose premise requires all three-state tests to fail? No. The condition
of the implication has not been met. Checking the premise is half of checking
a counterexample.

\clearpage
\section*{Lesson 2. Fractions, ratios, inequalities, and precision}
\addcontentsline{toc}{section}{2. Fractions, ratios, inequalities, and precision}
A fraction $a/b$, with $b\ne0$, means the number that becomes $a$ when multiplied
by $b$. Fractions are useful here because they can be exact. The decimal
$0.995$ equals $995/1000=199/200$ exactly. The displayed decimal $0.3333$ is
usually only an approximation to $1/3$. Knowing which situation applies matters
when a proof depends on the sign of a small difference.

\paragraph{A small exact calculation.}
The bath weights of the primary example are
\[
q_1=1,\qquad q_2=\frac{199}{200},\qquad
q_3=\frac25,\qquad q_4=\frac1{400}.
\]
The common denominator $400$ gives
\[
Q=q_1+q_2+q_3+q_4
=\frac{400+398+160+1}{400}=\frac{959}{400}.
\]
Dividing each weight by $Q$ therefore gives
\[
\left(\frac{400}{959},\frac{398}{959},
\frac{160}{959},\frac1{959}\right).
\]
You can verify normalization before touching a computer: the numerators sum to
$959$. Notice the distinction between a \emph{weight}, which is merely positive,
and a \emph{probability}, whose whole list must sum to one.

\paragraph{Ratios cancel shared scales.}
If $q_i$ and $q_j$ are both multiplied by the same positive number, $q_j/q_i$
does not change. Similarly, the ratio of two amplitudes is independent of a
shared normalization. This is why some intermediate choices are conventions
rather than physical assumptions. However, multiplying all energies by a
number without compensating the temperature generally changes the physics;
not every rescaling is harmless.

\paragraph{Inequalities carry direction.}
For $x<y$ and $c>0$, $cx<cy$. For $c<0$, $cx>cy$. A frequent mistake later is
to multiply an inequality by a negative quantity and forget to reverse it.
Another is to divide by a number before proving it is nonzero. The research
proofs pay particular attention to zero cases for this reason.

For example, $-0.094<-0.093$, but $|-0.094|>|-0.093|$. Absolute value measures
distance from zero and discards the sign. Our slow overlaps in the primary
cooling example are negative. A more negative overlap means a larger amplitude,
not a smaller one. We will compare absolute values explicitly.

\paragraph{Percentages need a denominator.}
Reducing a quantity from $100$ to $75$ is a $25\%$ reduction relative to $100$:
$(100-75)/100=1/4$. The reverse increase is $25/75=1/3$, not $25\%$. When we
define amplitude advantage as $1-|b_h|/|b_c|$, the denominator is the colder
preparation's amplitude. A different denominator would define a different
percentage.

\paragraph{Checkpoint.}
If $x\in[1.001,1.002]$ and $y\in[1.005,1.006]$, then $y-x\geq0.003>0$.
This proves an ordering even though neither value is known exactly. If the
intervals overlap, the same argument may be inconclusive. Inconclusive does
not mean false.

\clearpage
\section*{Lesson 3. A state is a list of probabilities}
\addcontentsline{toc}{section}{3. A state is a list of probabilities}
A finite-state model begins by choosing mutually exclusive alternatives. A
single system occupies one alternative at a time. The probability $p_i$ is the
chance of finding it in state $i$, or the limiting fraction of a large ensemble
found there. An ensemble means many independently prepared copies. The list
\[
p=(p_1,\ldots,p_n)^\top,\qquad p_i\geq0,
\qquad\sum_{i=1}^n p_i=1
\]
is a probability vector. The superscript $\top$ means ``transpose'': it turns a
horizontal list into a column. The symbol $\sum$ means add the indicated terms.
Thus $\sum_{i=1}^3p_i=p_1+p_2+p_3$.

\paragraph{Why a distribution contains more information than a mean.}
Take three energy levels $E=(0,1,2)$. The probability vectors
\[
p=(1/2,0,1/2)^\top,\qquad r=(0,1,0)^\top
\]
both have mean energy $1$. Nevertheless they are very different distributions:
one has no systems in the middle state, while the other has all systems there.
A statement about mean energy alone cannot in general determine a statement
about distance between full distributions. This distinction matters when
comparing different definitions of the Mpemba effect in the literature.

\paragraph{The geometry is simple before it is high-dimensional.}
For two states, $p_2=1-p_1$, so one number specifies the distribution. The
allowed points form a line segment. For three states, $p_3=1-p_1-p_2$ with
$p_1,p_2\geq0$ and $p_1+p_2\leq1$; the allowed points form a filled triangle.
For four states they form a solid tetrahedron. In general this set is called a
\emph{simplex}. The name is less important than its two rules: nonnegative
entries and total mass one.

\paragraph{Differences are not probabilities.}
If $p$ and $r$ are probability vectors, $z=p-r$ has entries summing to zero.
Its entries may be negative. For example,
\[
(0.6,0.4)^\top-(0.75,0.25)^\top=(-0.15,0.15)^\top.
\]
The negative first entry means a deficit relative to the reference, not a
negative physical probability. Eigenvectors describing deviations from
equilibrium likewise contain signs. Confusing a signed deviation with a
probability vector makes much of spectral theory seem paradoxical.

\paragraph{What the model leaves out.}
A finite-state model may describe energy levels, conformations, coarse regions
of a landscape, or another set of alternatives. Choosing these states is a
modeling decision. Mathematical consistency of their probability evolution
does not establish that they describe a particular experimental material.
That additional identification requires physical evidence.

\paragraph{Checkpoint.}
Is $(0.4,0.5,0.2)$ a probability vector? No: it sums to $1.1$. Is
$(-0.1,0.1,0)$ a possible difference of probability vectors? Yes. Its total is
zero, and appropriate reference probabilities make both original vectors
nonnegative.

\clearpage
\section*{Lesson 4. Averages, spread, and covariance}
\addcontentsline{toc}{section}{4. Averages, spread, and covariance}
An observable $f$ assigns a number $f_i$ to state $i$. Energy is one observable.
The mean or expectation of $f$ in distribution $p$ is
\[
\langle f\rangle_p=\sum_i p_i f_i.
\]
This is a weighted average, not an unweighted average unless all probabilities
are equal. For $p=(1/2,1/3,1/6)$ and $E=(0,1,2)$,
\[
\langle E\rangle_p=\tfrac12\,0+\tfrac13\,1+\tfrac16\,2=\tfrac23.
\]
The brackets $\langle\cdot\rangle$ are notation for an average; they do not
represent an extra operation beyond this sum.

\paragraph{Variance measures spread around the average.}
Define
\[
\Var_p(f)=\sum_i p_i(f_i-\langle f\rangle_p)^2
=\langle f^2\rangle_p-\langle f\rangle_p^2.
\]
To derive the second form, expand the square. The middle contribution is
$-2\langle f\rangle_p\sum_i p_i f_i=-2\langle f\rangle_p^2$, and the last
is $\langle f\rangle_p^2\sum_i p_i=\langle f\rangle_p^2$. Their sum gives
the stated identity. Because the first definition is a sum of nonnegative
terms, variance cannot be negative.

In our example, $\langle E^2\rangle=1/3+4/6=1$, so
$\Var(E)=1-(2/3)^2=5/9$. A zero variance would mean that the observable has
one value on every state with positive probability.

\paragraph{Covariance measures joint variation.}
For observables $f$ and $g$,
\[
\Cov_p(f,g)=\langle fg\rangle_p-\langle f\rangle_p\langle g\rangle_p.
\]
Positive covariance means that larger-than-average values of one tend, in the
weighted sum, to accompany larger-than-average values of the other. It is not
a causal statement. The same observable paired with itself gives variance.
A useful identity, obtained by expanding the products, is
\[
\Cov_p(f,g)=\frac12\sum_{i,j}p_ip_j(f_i-f_j)(g_i-g_j)
=\sum_{i<j}p_ip_j(f_j-f_i)(g_j-g_i).
\]
For the first equality, the two squared-index terms contribute
$2\langle fg\rangle$, and the two cross terms contribute
$-2\langle f\rangle\langle g\rangle$. The factor $1/2$ removes that two.
The second equality groups the equal $(i,j)$ and $(j,i)$ contributions.

\paragraph{Why this identity will matter.}
If $f$ and energy both increase with the state index, every pairwise product
$(f_j-f_i)(E_j-E_i)$ is nonnegative. Their covariance is then nonnegative.
No cancellation can change that sign. If $f$ is not ordered, different pairs
may compete. Later the observable $f$ will be a slow eigenfunction. Its shape
will determine whether varying the preparation temperature can reverse the
usual amplitude order.

\paragraph{Checkpoint.}
Adding a constant $c$ to every $f_i$ leaves $\Var(f)$ and $\Cov(f,g)$ unchanged.
This follows immediately from the difference form: $(f_j+c)-(f_i+c)=f_j-f_i$.
Scaling $f$ by $c$ scales its covariance by $c$ and its variance by $c^2$.

\clearpage
\section*{Lesson 5. Energy, temperature, and units}
\addcontentsline{toc}{section}{5. Energy, temperature, and units}
Energy levels are numbers $E_i$ assigned to the states. We label them in
increasing order,
\[
E_1<E_2<\cdots<E_n.
\]
The temperature $T$ is positive. Statistical mechanics often uses inverse
temperature
\[
\beta=\frac{1}{k_B T},
\]
where $k_B$ is Boltzmann's constant. The product $\beta E_i$ is dimensionless.
An exponential of an energy measured in joules would be meaningless without
such a conversion; the number placed in an exponential must carry no unit.

\paragraph{The inverse reverses the ordering.}
Higher temperature means lower $\beta$. With a colder bath at $T_b$, two
cooling preparations satisfy
\[
T_h>T_c>T_b
\quad\Longleftrightarrow\quad
0<\beta_h<\beta_c<\beta_b.
\]
The subscripts $h$ and $c$ mean hotter and colder \emph{relative to each other}.
Both are hotter than the bath. In a heating experiment they can both be colder
than the bath. It is worth writing the inverse-temperature ordering on paper
before every sign argument.

\paragraph{Changing the energy zero changes no probabilities.}
Only energy differences matter for the Gibbs probabilities introduced next.
If $E_i$ is replaced by $E_i+C$, every exponential weight is multiplied by
$e^{-\beta C}$. That common factor cancels in normalization. We can therefore
choose $E_1=0$. This choice does not say that the ground state has no physical
energy; it selects a reference level.

\paragraph{Choosing bath units.}
Let $\widetilde E_i=\beta_b(E_i-E_1)$ and
$\widetilde\beta=\beta/\beta_b$. Then
$\widetilde\beta\widetilde E_i=\beta(E_i-E_1)$, so the probabilities are
unchanged. In these units the bath inverse temperature is $1$. A preparation
with $\widetilde\beta=1/10$ is at ten times the bath's absolute temperature.
Dropping the tildes simplifies notation, but does not turn inverse temperature
into ordinary temperature.

\paragraph{Do not confuse three rescalings.}
Shifting all energies changes their zero. Scaling energy and inverse
temperature reciprocally changes their units. Multiplying every transition
rate by a positive factor changes the clock. These are distinct operations.
Setting an energy gap to one and also setting the bath inverse temperature to
one cannot generally be done as two independent harmless choices; their
product contains physical information.

\paragraph{Checkpoint.}
If the bath has $\beta_b=1$, then $\beta=0.01$ is a hot preparation, not a cold
one. The limiting value $\beta=0$ represents infinite temperature. A proof at
that endpoint alone is not yet a proof for two finite initial temperatures;
a continuity argument or finite choices must complete it.

\clearpage
\section*{Lesson 6. Why exponentials and logarithms appear}
\addcontentsline{toc}{section}{6. Why exponentials and logarithms appear}
The exponential $e^x$ is the positive function with $e^0=1$ and derivative
$e^x$. It obeys $e^{x+y}=e^xe^y$. One definition that also supports rigorous
computation is
\[
e^x=1+x+\frac{x^2}{2!}+\frac{x^3}{3!}+\cdots,
\]
where $m!=1\cdot2\cdots m$, with $0!=1$. For a real number $x$, the series
converges. For negative $x$, $e^x=1/e^{-x}$ is still positive. It never becomes
negative just because the exponent is negative.

\paragraph{Repeated fractional change.}
A quantity that loses a fixed \emph{fraction per unit time} satisfies
$y'(t)=-k y(t)$ with $k>0$. The prime means the instantaneous rate of change; Lesson 8 defines it
precisely. Its solution is $y(t)=y(0)e^{-kt}$. Differentiating
checks the equation directly. At time $1/k$ the quantity is multiplied by
$e^{-1}$, about $0.368$. The number $1/k$ is a characteristic decay time, not
the time at which the quantity becomes exactly zero. An exponential with a
nonzero initial value reaches zero only in the limit $t\to\infty$.

\paragraph{Logarithms undo exponentials.}
For $a>0$, $\log a$ is the real number satisfying $e^{\log a}=a$.
We use the natural logarithm throughout. Thus $\log(ab)=\log a+\log b$ and
$\log(a/b)=\log a-\log b$. Solving $q=e^{-E}$ gives $E=-\log q$.
For $0<q<1$, $\log q<0$, so this energy is positive. Smaller weights correspond
to larger energies.

The primary bath weights therefore give
\[
E_1=0,\quad E_2=\log(200/199),\quad
E_3=\log(5/2),\quad E_4=\log400.
\]
Their order follows without decimal approximations because
$1<200/199<5/2<400$ and the logarithm is increasing.

\paragraph{Real powers connect preparation to bath.}
With $q_i=e^{-E_i}$ in bath units,
\[
e^{-\beta E_i}=e^{\beta\log q_i}=q_i^\beta.
\]
For $\beta=1/2$, this is a square root. For $\beta=1/1000$, it is a positive
thousandth root. A rational bath weight need not produce a rational hot
preparation weight. That does not make exact certification impossible:
rational intervals can enclose the positive root, with their validity checked
by integer powers.

\paragraph{Checkpoint.}
The equation $e^{-kt}=a$ with $0<a<1$ gives
$t=-\log(a)/k>0$. Taking logarithms converts a multiplicative decay target into
an additive time formula. This simple fact will later explain why an amplitude
percentage and a time-saving percentage are different.

\clearpage
\section*{Lesson 7. Constructing a Gibbs distribution by hand}
\addcontentsline{toc}{section}{7. Constructing a Gibbs distribution by hand}
The thermal or Gibbs preparation at inverse temperature $\beta$ is
\[
\pi_i(\beta)=\frac{e^{-\beta E_i}}{Z(\beta)},
\qquad Z(\beta)=\sum_j e^{-\beta E_j}.
\]
The symbol $Z$ is called the partition function. At this stage, its job is
simply to normalize positive weights. Every numerator is positive and the
numerators sum to $Z$, so the probabilities are positive and sum to one.

\paragraph{Illustration with three states.}
Use $E=(0,\log2,\log4)$ and $\beta=1$. The unnormalized weights are
$(1,1/2,1/4)$, because $e^{-\log2}=1/2$. Their sum is $7/4$. Dividing gives
\[
\pi(1)=(4/7,2/7,1/7)^\top.
\]
At $\beta=0$, every weight is $1$, so $\pi(0)=(1/3,1/3,1/3)^\top$.
At $\beta=2$, the weights are squared, giving $(1,1/4,1/16)$ and hence
$\pi(2)=(16/21,4/21,1/21)^\top$. Increasing inverse temperature concentrates
probability in the low-energy states.

\paragraph{Ratios are even simpler than probabilities.}
Normalization cancels from a probability ratio:
\[
\frac{\pi_j(\beta)}{\pi_i(\beta)}
=e^{-\beta(E_j-E_i)}.
\]
For the example, the ratio of the third probability to the first at $\beta=1$
is $1/4$. At $\beta=2$ it is $1/16$. This is how energy gaps control relative
occupancy. The value of a single probability depends on every state through
$Z$, while the ratio of two probabilities depends only on their energy gap.

\paragraph{The bath and the preparation have different jobs.}
We write $\pi_b=\pi(\beta_b)$ for the bath equilibrium. This fixed distribution
helps determine the dynamics. We vary $\beta$ to choose initial distributions
$p(0)=\pi(\beta)$. Once the bath is fixed, varying the preparation does not
change the matrix of transition rates. Comparing two preparations with two
different matrices would be a different experiment.

\paragraph{What happens after time zero?}
Although $p(0)$ lies on the Gibbs curve $\beta\mapsto\pi(\beta)$, the evolving
probability vector usually leaves that curve. One cannot generally assign a
single instantaneous temperature to the entire vector. The master equation
must evolve all probabilities, rather than just a scalar temperature.

\paragraph{Checkpoint.}
In the illustration,
$\langle E\rangle_{\pi(1)}=(4/7)\log2$, and
$\Var_{\pi(1)}(E)=(26/49)(\log2)^2$. To check the variance, first compute
$\langle E^2\rangle=(6/7)(\log2)^2$ and subtract the square of the mean.
The three-state illustration is a teaching example, not the four-state
counterexample.

\clearpage
\section*{Lesson 8. Derivatives, local change, and Taylor expansion}
\addcontentsline{toc}{section}{8. Derivatives, local change, and Taylor expansion}
For a function $f(x)$, the derivative is the limiting slope
\[
f'(x)=\lim_{h\to0}\frac{f(x+h)-f(x)}{h}.
\]
A positive derivative means increasing values locally; a negative derivative
means decreasing values locally. A zero derivative does not alone identify a
maximum. The function $x^3$ has derivative zero at $0$ but keeps increasing
through it. The function $x^2$ has a minimum there. One needs signs on either
side or additional information about higher derivatives.

\paragraph{The rules we use most often.}
The product rule is $(fg)'=f'g+fg'$. The quotient rule is
$(f/g)'=(f'g-fg')/g^2$ when $g\ne0$. The chain rule gives
$(e^{u(x)})'=u'(x)e^{u(x)}$. These rules make the Gibbs derivative accessible.
Let $r_i=e^{-\beta E_i}$. Then $r_i'=-E_i r_i$ and
\[
Z'=\sum_i r_i'=-\sum_i E_i r_i=-Z\langle E\rangle_\pi.
\]
Therefore
\begin{align*}
\pi_i'&=\frac{r_i'Z-r_iZ'}{Z^2}\\
&=\frac{-E_i r_i Z+r_i Z\langle E\rangle_\pi}{Z^2}\\
&=-(E_i-\langle E\rangle_\pi)\pi_i.
\end{align*}
Every line follows from substitution and cancellation. A state below the
current mean energy gains probability as $\beta$ increases; a state above it
loses probability. Summing the last line gives zero, as required because the
total probability remains one.

\paragraph{Taylor expansion is a controlled local approximation.}
Near a point $x$, a sufficiently smooth function has
\[
f(x+h)=f(x)+h f'(x)+\tfrac12h^2f''(x)+\text{remainder}.
\]
An approximation without a remainder estimate is useful intuition, but not a
rigorous numerical enclosure. For example,
$\log(1+u)=u-u^2/2+O(u^3)$ as $u\to0$. Here $O(u^3)$ means that the absolute
remainder is bounded by a constant times $|u|^3$ sufficiently near zero.
It does not mean the remainder is literally zero.

The notation $r(h)=o(h^2)$ is stronger than merely tending to zero: it means
$r(h)/h^2\to0$. These symbols will let us separate the leading late-time term
from faster-decaying errors. Their constants and domains matter whenever an
asymptotic statement is turned into a finite numerical guarantee.

\paragraph{Checkpoint.}
At the three-state Gibbs distribution from Lesson 7,
$\pi_1'=\langle E\rangle\pi_1=(16/49)\log2>0$. The ground-state population
increases as the system is prepared colder. The derivative is with respect to
inverse temperature, not time.

\clearpage
\section*{Lesson 9. The covariance identity that drives the mechanism}
\addcontentsline{toc}{section}{9. The covariance identity that drives the mechanism}
Keep an observable $f_i$ fixed while varying the Gibbs preparation. Using
Lesson 8,
\begin{align*}
\frac{d}{d\beta}\langle f\rangle_\pi
&=\sum_i f_i\pi_i'\\
&=-\sum_i f_i E_i\pi_i+\langle E\rangle_\pi\sum_i f_i\pi_i\\
&=-\langle fE\rangle_\pi+\langle E\rangle_\pi\langle f\rangle_\pi\\
&=-\Cov_\pi(f,E).
\end{align*}
This identity translates a temperature derivative into a weighted correlation.
It is exact, not a numerical approximation. The word ``fixed'' is essential:
if $f$ itself depended on $\beta$, an additional term
$\sum_i\pi_i f_i'$ would appear.

\paragraph{First application: mean energy.}
Choose $f=E$. Then
\[
\frac{d}{d\beta}\langle E\rangle_\pi=-\Var_\pi(E)<0
\]
for distinct energies and a finite Gibbs distribution. The strict sign follows
because all probabilities are positive and not all energies are equal.
Thus mean energy decreases when inverse temperature increases. This is the
expected behavior of thermal preparations, derived directly from the formula.

\paragraph{Second application: a shaped observable.}
If $f_1\leq f_2\leq\cdots\leq f_n$, the pairwise covariance formula makes
$\Cov(f,E)\geq0$. Its thermal average cannot increase with $\beta$.
If the values of $f$ bend up and down instead, pair contributions can have
opposite signs. Temperature changes their weights. Their competition can
therefore change the sign of the derivative.

This is the basic mathematical opening for an anomalous slow amplitude. The
observable that matters is not simply energy; it is an eigenfunction of the
dynamics. A hot preparation can have more average energy and nevertheless a
smaller projection onto that particular eigenfunction. There is no
contradiction because these are different weighted sums.

\paragraph{How to use the identity without overusing it.}
A positive covariance at one preparation gives a negative derivative there.
It does not automatically describe every temperature. A zero covariance marks
a stationary point of the average, but the derivative may touch zero without
changing sign. To prove a reversal, we need either a sign change with the
right amplitude orientation or two strictly ordered amplitude values.
The exact counterexample uses the latter route.

\paragraph{Checkpoint.}
The bath determines the eigenfunction. Holding that bath and the generator
fixed is precisely what allows the covariance identity to be used for the
slow overlap. Differentiating through a family of changing baths would be a
different calculation and would need the missing derivative terms.

\clearpage
\section*{Lesson 10. Transition rates and a complete two-state solution}
\addcontentsline{toc}{section}{10. Transition rates and a complete two-state solution}
A transition rate is a probability per unit time in the small-time limit.
If the system is in state $j$, a rate $a_{ij}$ means that the probability of
jumping to a different state $i$ during a short interval $h$ is
$a_{ij}h+o(h)$. The rate itself need not be less than one: it has units of
inverse time, whereas the small-time jump probability is dimensionless.
A rate of $3$ per second is perfectly legitimate.

\paragraph{Destination comes first.}
Throughout this report $a_{ij}$ means a jump \emph{from $j$ to $i$}.
For ordered energies, $a_{12}$ is a downward transition from state $2$ to
state $1$. This convention is easy to reverse accidentally. Read a matrix
column as the possible destinations from its starting state.

\paragraph{A model we can solve without matrix machinery.}
Let the upward rate from state $1$ to $2$ be $1$ and the downward rate from
$2$ to $1$ be $3$. Incoming flow to state $2$ is $1p_1$; outgoing flow is
$3p_2$. Thus
\[
p_2'=p_1-3p_2=(1-p_2)-3p_2=1-4p_2.
\]
The equilibrium is found by setting the derivative to zero:
$p_{2b}=1/4$ and $p_{1b}=3/4$. Define $z=p_2-1/4$. Its derivative is
$z'=p_2'=1-4(z+1/4)=-4z$, so Lesson 6 gives
\[
p_2(t)=\frac14+\left(p_2(0)-\frac14\right)e^{-4t},
\qquad p_1(t)=1-p_2(t).
\]
This is a complete solution for every initial probability vector.

\paragraph{Check the solution rather than trusting it.}
At $t=0$ it gives the chosen initial value. Differentiating yields
$-4(p_2(0)-1/4)e^{-4t}=1-4p_2(t)$, so it satisfies the equation.
As $t\to\infty$, the exponential goes to zero and $p_2(t)\to1/4$.
For $t\geq0$, the solution is a convex combination of $p_2(0)$ and $1/4$,
meaning a weighted average with nonnegative weights that sum to one,
so it remains between zero and one. Conservation, positivity, and relaxation
have each been checked separately.

\paragraph{Why this model does not exhibit the desired reversal.}
For thermal preparations above this bath, the hotter preparation has the
larger state-$2$ population. Its initial deviation from $1/4$ is therefore
larger. Both deviations are multiplied by the same positive $e^{-4t}$.
Their ordering cannot reverse. A single scalar decay leaves no room for
competition among differently decaying components. Extra states create
additional components, although extra states alone do not guarantee a reversal.

\clearpage
\section*{Lesson 11. Matrices are bookkeeping for many flows}
\addcontentsline{toc}{section}{11. Matrices are bookkeeping for many flows}
A matrix is a rectangular array of numbers. Multiplying a matrix by a column
vector means taking one row at a time and adding entrywise products.
Our two-state model becomes
\[
\frac{d}{dt}\begin{pmatrix}p_1\\p_2\end{pmatrix}
=\underbrace{\begin{pmatrix}-1&3\\1&-3\end{pmatrix}}_{M}
\begin{pmatrix}p_1\\p_2\end{pmatrix}.
\]
The first row says $p_1'=-p_1+3p_2$; the second says $p_2'=p_1-3p_2$.
Nothing new has been assumed by writing the two equations compactly.

\paragraph{The master equation for $n$ states.}
Set $M_{ij}=a_{ij}\geq0$ for $i\ne j$, and put the negative total exit rate
on the diagonal:
\[
M_{jj}=-\sum_{i\ne j}a_{ij}.
\]
Then $p'=Mp$ means
\[
p_i'=\sum_{j\ne i}a_{ij}p_j
-\left(\sum_{j\ne i}a_{ji}\right)p_i.
\]
The first sum is incoming probability per time, and the second is outgoing
probability per time. Notice the reversed indices in the outgoing rate.
A diagonal entry is not a negative jump probability. It records the loss
from that state.

\paragraph{Conservation comes from columns.}
Every column of $M$ sums to zero. Therefore
\[
\frac{d}{dt}\sum_i p_i
=\sum_i\sum_jM_{ij}p_j
=\sum_j\left(\sum_iM_{ij}\right)p_j=0.
\]
If the initial total is one, it stays one. Row sums need not vanish under
this column-vector convention. Accidentally imposing row conservation
would describe a transposed convention rather than the stated model.

\paragraph{Why nonnegative probabilities remain nonnegative.}
Choose $r$ at least as large as every exit rate, with $r>0$, and define
$P=I+M/r$. Here $I$ is the identity matrix, with ones on its diagonal and
zeros elsewhere. Every entry of $P$ is nonnegative and every column sums
to one. Thus $P$ maps probability vectors to probability vectors. The exact
solution can be written
\[
e^{Mt}=e^{-rt}\sum_{k=0}^{\infty}\frac{(rt)^k}{k!}P^k.
\]
The coefficients are nonnegative and sum to one. This expresses the solution
as a weighted average of probability-preserving maps, establishing positivity.
The matrix exponential is defined by the same power series as the scalar
exponential, with matrix powers replacing scalar powers.

\paragraph{Checkpoint.}
An Euler step $p(t+h)\approx(I+hM)p(t)$ can give negative probabilities if
$h$ is too large. That is a failure of the numerical step, not of the exact
continuous-time process. The series above also becomes a useful certificate
method because its omitted tail can be bounded explicitly.

\clearpage
\section*{Lesson 12. Equilibrium, stationarity, and detailed balance}
\addcontentsline{toc}{section}{12. Equilibrium, stationarity, and detailed balance}
A stationary distribution $\pi_b$ satisfies $M\pi_b=0$. Starting there gives
$p(t)=\pi_b$ for all time. Stationarity says that total incoming and outgoing
flow cancel at each state. A stronger condition, detailed balance, says that
each pair balances separately:
\[
M_{ij}\pi_{bj}=M_{ji}\pi_{bi}\qquad(i\ne j).
\]
The two sides are equilibrium probability currents in opposite directions.
Their equality makes the dynamics reversible in the probabilistic sense.
It does not mean that a nonequilibrium preparation fails to relax.

\paragraph{Detailed balance implies stationarity.}
For a fixed state $i$,
\begin{align*}
(M\pi_b)_i
&=\sum_{j\ne i}M_{ij}\pi_{bj}+M_{ii}\pi_{bi}\\
&=\sum_{j\ne i}M_{ji}\pi_{bi}
-\sum_{j\ne i}M_{ji}\pi_{bi}=0.
\end{align*}
The first equality uses matrix multiplication; the second uses detailed
balance and the diagonal definition. Every cancellation is visible.

\paragraph{Constructing reversible rates from weights.}
Since $\pi_{bi}=q_i/Q$, detailed balance is equivalent to
$a_{ij}q_j=a_{ji}q_i$. Choose a downward rate $a_{ij}>0$ for $i<j$ and set
\[
a_{ji}=a_{ij}\frac{q_j}{q_i}.
\]
All weights are positive, so the upward rate is positive. If $q_j<q_i$,
the upward rate is smaller than its downward partner. Choosing the diagonal
last ensures conservation. This construction makes the desired bath
stationary automatically.

\paragraph{An explicit check.}
For the two-state example, take $q=(1,1/3)$ and downward rate $a_{12}=3$.
Then $a_{21}=3(1/3)=1$, as before. Normalization gives
$\pi_b=(3/4,1/4)$. The flow from $2$ to $1$ is $3(1/4)=3/4$;
the flow from $1$ to $2$ is $1(3/4)=3/4$. Individual transitions continue
at equilibrium even though the probability distribution does not change.

\paragraph{Connectivity matters.}
If every off-diagonal rate is positive, any state can reach any other.
There is then one stationary distribution and relaxation to it. More
sparse models can also have this property when their transition graph is
connected appropriately. If the graph splits into disconnected components,
there can be several stationary directions. Some proof steps survive that
boundary and others require separate treatment; positivity assumptions are
therefore more than decorative wording.

\clearpage
\section*{Lesson 13. Eigenvectors separate the decay patterns}
\addcontentsline{toc}{section}{13. Eigenvectors separate the decay patterns}
An eigenvector $v\ne0$ of $M$ satisfies $Mv=\lambda v$. The matrix changes
it by the factor $\lambda$ while keeping it on the same line. A negative
factor reverses its sense along that line. The number
$\lambda$ is an eigenvalue. For the differential equation, a component in
that direction evolves as $e^{\lambda t}v$, because
\[
\frac{d}{dt}(e^{\lambda t}v)=\lambda e^{\lambda t}v
=M(e^{\lambda t}v).
\]
An eigenvalue is therefore a decay rate in the dynamics, not a probability.

\paragraph{The two-state eigenvectors can be checked by multiplication.}
For $M=\left(\begin{smallmatrix}-1&3\\1&-3\end{smallmatrix}\right)$,
\[
M\begin{pmatrix}3/4\\1/4\end{pmatrix}=0,
\qquad
M\begin{pmatrix}1\\-1\end{pmatrix}
=\begin{pmatrix}-4\\4\end{pmatrix}
=-4\begin{pmatrix}1\\-1\end{pmatrix}.
\]
The stationary mode has eigenvalue zero. The deviation mode has eigenvalue
$-4$. Its entries sum to zero, so adding a small multiple to the equilibrium
preserves total probability.

\paragraph{Several modes give several clocks.}
For a reversible connected four-state model, the eigenvalues can be ordered
\[
0=\lambda_1>\lambda_2\geq\lambda_3\geq\lambda_4.
\]
The nonzero values are negative. If $\lambda_2>\lambda_3$, the slowest
nonstationary mode is simple, meaning one-dimensional. Its exponential
decays less rapidly than the others: $e^{-t}$ outlasts $e^{-10t}$.
At late times, even a small coefficient of the slower exponential can dominate.

\paragraph{A two-exponential illustration.}
Consider the positive functions
$f_h(t)=e^{-t}+9e^{-10t}$ and $f_c(t)=2e^{-t}+e^{-10t}$.
At $t=0$, $f_h=10>3=f_c$. Their difference is
$-e^{-t}+8e^{-10t}$, which is zero exactly when $e^{9t}=8$.
Thus they cross at $t=(\log8)/9$, and $f_h<f_c$ thereafter.
This is an illustration of competing decay coefficients, not a complete
thermal Markov counterexample. Its role is to make the mechanism imaginable:
more initial total deviation can coexist with less slow deviation.

\paragraph{Checkpoint.}
The slow eigenvalue is the negative one closest to zero, not the most negative
one. A very negative eigenvalue gives a very fast decay. A simple slow mode
permits one scalar amplitude; a repeated slow eigenvalue requires a whole
projection, discussed in Lesson 26.

\clearpage
\section*{Lesson 14. Left eigenvectors measure the mode amplitudes}
\addcontentsline{toc}{section}{14. Left eigenvectors measure the mode amplitudes}
A right eigenvector describes a direction of change in the probability
vector. A left eigenvector is a row vector satisfying
$w^\top M=\lambda w^\top$, equivalently $M^\top w=\lambda w$.
Multiplying the state by this row gives a scalar mode coordinate. Indeed,
\[
\frac{d}{dt}(w^\top p)=w^\top Mp=\lambda w^\top p.
\]
Thus $w^\top p(t)=e^{\lambda t}w^\top p(0)$ exactly.

\paragraph{Why the bath disappears from the overlap.}
For a nonzero eigenvalue,
\[
0=w^\top M\pi_b=\lambda w^\top\pi_b.
\]
Dividing by $\lambda\ne0$ gives $w^\top\pi_b=0$. Consequently the
coefficient can be computed as either $w^\top[p(0)-\pi_b]$ or
$w^\top p(0)$. For a Gibbs preparation we call it
\[
b(\beta)=w^\top\pi(\beta)=\sum_i w_i\pi_i(\beta).
\]
It is a weighted average of the values of the eigenfunction $w$.

\paragraph{Normalization prevents a hidden factor.}
Choose a corresponding right eigenvector $v$ such that $w^\top v=1$.
Then $b(\beta)v$ is the actual vector component in that mode. If $w$ is
multiplied by $c\ne0$, $b$ is multiplied by $c$ and $v$ must be divided by
$c$. The physical product does not change. Ratios $|b_h|/|b_c|$ are therefore
independent of this common normalization.

\paragraph{Return to the complete two-state solution.}
Take $w=(1,-3)^\top$. Direct multiplication gives $w^\top M=-4w^\top$.
With $v=(1,-1)^\top/4$, $w^\top v=(1+3)/4=1$. The initial overlap is
$b=p_1-3p_2=1-4p_2$. Therefore
\[
p(t)=\pi_b+e^{-4t}(1-4p_2(0))\frac{(1,-1)^\top}{4}.
\]
Its second component equals the scalar formula of Lesson 10. The abstract
left/right distinction has recovered the elementary answer exactly.

\paragraph{Covariance now acquires its physical role.}
Because $w$ is fixed when the preparation varies,
\[
b'(\beta)=-\Cov_{\pi(\beta)}(w,E).
\]
This is the formula derived in Lesson 9 with $f=w$. The thermal curve and the
dynamical mode meet in a single weighted covariance. The rates choose $w$;
the energies and preparation temperature choose the weights in the covariance.
The effect depends on both pieces.

\clearpage
\section*{Lesson 15. Reversibility turns a difficult matrix into a symmetric one}
\addcontentsline{toc}{section}{15. Reversibility and the symmetric representation}
A real symmetric matrix equals its transpose. Such matrices have real
eigenvalues and an orthonormal basis of eigenvectors: mutually perpendicular
unit-length vectors in which every vector has a unique expansion. Our probability
generator is generally not symmetric in the ordinary coordinates, because
upward and downward rates differ. Detailed balance provides the right change
of coordinates.

Let $D=\diag(\pi_b)$ be the diagonal matrix of bath probabilities. Its square
root has entries $\sqrt{\pi_{bi}}$, and its inverse square root has entries
$1/\sqrt{\pi_{bi}}$. All are defined because the probabilities are positive.
Set
\[
S=D^{-1/2}MD^{1/2}.
\]
Detailed balance says $MD=DM^\top$, and therefore
\[
S^\top=D^{1/2}M^\top D^{-1/2}
=D^{-1/2}MD^{1/2}=S.
\]
This change of coordinates is called a similarity transformation. It leaves
eigenvalues unchanged: if $Mv=\lambda v$, then
$S(D^{-1/2}v)=\lambda(D^{-1/2}v)$.

\paragraph{Why decay rates cannot be positive.}
Define equilibrium conductances $c_{ij}=M_{ij}\pi_{bj}=c_{ji}$ for $i\ne j$.
Writing the generator on observables as $L=M^\top$, direct expansion gives
\[
-f^\top DLf=\sum_{i<j}c_{ij}(f_i-f_j)^2\geq0.
\]
For an eigenfunction $Lf=\lambda f$, the left side equals
$-\lambda\sum_i\pi_{bi}f_i^2$. The weighted squared length is positive
when $f\ne0$, so $\lambda\leq0$. If all conductances are positive, equality
requires all values $f_i$ to coincide. Only the constant observable has
zero eigenvalue, corresponding to the unique stationary probability mode.

\paragraph{The right inner product.}
For observables use $\langle f,g\rangle_D=\sum_i\pi_{bi}f_i g_i$.
For probability deviations use
$\langle z,y\rangle_{D^{-1}}=\sum_i z_i y_i/\pi_{bi}$.
These weighted inner products are the natural ones for detailed balance.
The different weights arise because a probability deviation and an
observable are dual objects. If $Lw=\lambda w$, then $Dw$ is a right
eigenvector of $M$.

\paragraph{Advanced consequence, in plain terms.}
The spectral analysis is not relying on a lucky numerical property of one
matrix. Reversibility supplies a real diagonalization throughout the stated
domain. That removes oscillatory complex modes and nontrivial Jordan blocks
from this problem. It does not remove the need to check whether the slow
eigenvalue is simple or repeated.

\clearpage
\section*{Lesson 16. What does it mean to be closer to equilibrium?}
\addcontentsline{toc}{section}{16. What does it mean to be closer to equilibrium?}
To compare whole probability vectors we need a specified distance or
divergence. Different choices emphasize different discrepancies. Three used
in this project are
\begin{align*}
\mathrm{TV}(p,\pi_b)&=\frac12\sum_i|p_i-\pi_{bi}|,\\
\chi^2(p\Vert\pi_b)&=\sum_i\frac{(p_i-\pi_{bi})^2}{\pi_{bi}},\\
D_{\mathrm{KL}}(p\Vert\pi_b)&=\sum_i p_i\log\frac{p_i}{\pi_{bi}}.
\end{align*}
A norm assigns a length to a vector: it is positive except at zero,
scales by the absolute value of a scalar multiplier, and obeys the triangle
inequality $\|x+y\|\leq\|x\|+\|y\|$. Euclidean length
$\sqrt{\sum_i x_i^2}$ is one example. A metric measures distance between
two points and also requires symmetry and the triangle inequality.
Total variation (TV) is a metric. Pearson chi-square and relative entropy
(KL) are divergences; they need not satisfy every metric axiom. We use the
word ``distance'' informally only when the specified formula is clear.
For KL, a zero $p_i$ contributes zero by a limit. The bath entries are positive.

\paragraph{A hand calculation.}
Take $\pi_b=(3/4,1/4)$ and $p=(1/2,1/2)$. The deviations are
$(-1/4,1/4)$, so
\[
\mathrm{TV}=\tfrac12(\tfrac14+\tfrac14)=\tfrac14,
\qquad
\chi^2=\frac{1/16}{3/4}+\frac{1/16}{1/4}=\tfrac13.
\]
The KL value is
$\frac12\log(2/3)+\frac12\log2=\frac12\log(4/3)$.
The numerical values differ because the measurements differ. Comparing a TV
number with a KL number as though they used the same scale would be meaningless.

\paragraph{Why KL is approximately quadratic near equilibrium.}
Write $p_i=\pi_{bi}(1+u_i)$, where $\sum_i\pi_{bi}u_i=0$.
For small $u$,
\[
(1+u)\log(1+u)=u+\tfrac12u^2+O(u^3).
\]
Multiplying by $\pi_{bi}$ and summing cancels the linear term. Hence
\[
D_{\mathrm{KL}}(p\Vert\pi_b)
=\tfrac12\sum_i\pi_{bi}u_i^2+O\!\left(\sum_i\pi_{bi}|u_i|^3\right)
=\tfrac12\chi^2(p\Vert\pi_b)+\text{higher-order terms}.
\]
This is a local expansion around the bath, not an identity at arbitrary
probabilities. Rare bath states have large reciprocal weights in chi-square,
which is one reason finite-time comparisons require care.

\paragraph{Initially hotter is initially farther in KL.}
For a Gibbs preparation,
\[
D_{\mathrm{KL}}(\pi(\beta)\Vert\pi_b)
=(\beta_b-\beta)\langle E\rangle_{\pi(\beta)}
+\log Z(\beta_b)-\log Z(\beta).
\]
Differentiate, use $\langle E\rangle'=-\Var(E)$ and
$(\log Z)'=-\langle E\rangle$, and the two mean-energy terms cancel:
\[
\frac{d}{d\beta}D_{\mathrm{KL}}=(\beta-\beta_b)\Var(E).
\]
On the cooling side this derivative is negative. Thus increasing $\beta$
toward the bath decreases the initial KL divergence, and the hotter
preparation starts farther away. This supplies the initial half of a genuine
initial-to-eventual reversal.

\clearpage
\section*{Lesson 17. Turning a slow-overlap inequality into a reversal}
\addcontentsline{toc}{section}{17. From a slow-overlap inequality to a reversal}
Assume the slow nonzero eigenvalue is simple, so
$\lambda_2>\lambda_3$. Decompose the initial deviation into eigenmodes.
At late times,
\[
p(t;\beta)-\pi_b=e^{\lambda_2t}b(\beta)v_2+R(t;\beta),
\qquad R(t;\beta)=O(e^{\lambda_3t}).
\]
After division by $e^{\lambda_2t}$, the remainder tends to zero because
$e^{(\lambda_3-\lambda_2)t}\to0$. This is why the slow coefficient determines
the eventual order.

For any fixed norm, continuity and homogeneity give
\[
\lim_{t\to\infty}e^{-\lambda_2t}\|p(t;\beta)-\pi_b\|
=|b(\beta)|\|v_2\|.
\]
If $|b_h|<|b_c|$, the two limiting values have a strictly positive gap.
For sufficiently late finite times, the errors are smaller than a fraction
of that gap, so the same strict ordering holds. This is an existence argument
for a finite late-time interval, not merely an assertion at $t=\infty$.

\paragraph{For KL the coefficient is squared.}
Combining the spectral expansion with Lesson 16 gives
\[
D_{\mathrm{KL}}(p(t;\beta)\Vert\pi_b)
=\frac12e^{2\lambda_2t}b(\beta)^2
\sum_i\frac{v_{2i}^2}{\pi_{bi}}+o(e^{2\lambda_2t}).
\]
The positive factor involving $v_2$ is shared by both preparations.
A smaller absolute overlap therefore gives smaller eventual KL divergence.
Since the hotter Gibbs preparation starts farther in KL, continuity forces
at least one crossing between the initial and eventual orderings.

\paragraph{Ordinary and strong effects.}
An ordinary effect can occur with both overlaps nonzero. A strong effect
requires $b(\beta)=0$ at a preparation different from the bath, eliminating
the slow mode altogether. The bath itself always has zero overlap and does
not count as a nontrivial strong preparation. The primary counterexample
uses two nonzero overlaps; it does not need a strong effect.

\paragraph{Amplitude advantage is not time advantage.}
Define $A=1-|b_h|/|b_c|$. If $A=1/4$, the slow amplitude ratio is $3/4$,
and the asymptotic KL ratio is $9/16$. Neither number is automatically a
$25\%$ reduction in target time. In a pure one-mode KL approximation with
$\lambda_2=-r<0$, the time to a small threshold $d$ is
\[
t(\beta;d)\approx\frac{1}{2r}
\log\frac{K b(\beta)^2}{d},
\qquad
 t_c-t_h\approx\frac1r\log\frac{|b_c|}{|b_h|}.
\]
The difference is approximately additive in time. Its fraction of $t_c$
depends on the threshold and tends to zero as the threshold tends to zero
with fixed nonzero overlaps. Finite-time claims require their own proof.

\clearpage
\section*{Lesson 18. Assemble the exact four-state witness}
\addcontentsline{toc}{section}{18. Assemble the exact four-state witness}
We now have the ingredients to read actual research data. The primary
witness uses the weights from Lesson 2, bath inverse temperature $1$, and
six downward rates
\[
(a_{12},a_{13},a_{14},a_{23},a_{24},a_{34})
=\left(1,28,\frac{249}{10},\frac{1799}{8},
\frac{4983}{200},\frac{251}{25}\right).
\]
These are inputs, not rounded numerical estimates. The reverse rates are
computed from detailed balance, and the diagonal entries from conservation.

\paragraph{Build one column completely.}
The three exits from state $1$ are upward transitions:
\begin{align*}
a_{21}&=a_{12}\frac{q_2}{q_1}=\frac{199}{200},\\
a_{31}&=a_{13}\frac{q_3}{q_1}=28\frac25=\frac{56}{5},\\
a_{41}&=a_{14}\frac{q_4}{q_1}=\frac{249}{10}\frac1{400}
=\frac{249}{4000}.
\end{align*}
Their sum is $(3980+44800+249)/4000=49029/4000$.
The first column is therefore
\[
\left(-\frac{49029}{4000},\frac{199}{200},
\frac{56}{5},\frac{249}{4000}\right)^\top.
\]
It sums to zero exactly. Repeating this operation for the other states gives
the matrix printed in Part II. You do not have to infer its meaning from a
wall of fractions: every positive entry is a jump rate, and every negative
diagonal entry is the negative sum of exits.

\paragraph{Check one less obvious reverse rate.}
For the pair $2,3$,
\[
a_{32}=\frac{1799}{8}\frac{2/5}{199/200}
=\frac{1799}{8}\frac{80}{199}=\frac{17990}{199}.
\]
It follows that $a_{23}q_3=a_{32}q_2$ exactly. Such checks expose a reversed
ratio immediately; using $q_2/q_3$ would produce the wrong bath.

\paragraph{The two preparations.}
Take $\beta_h=1/1000$ and $\beta_c=1/100$. Both are positive and smaller
than the bath value $1$, so both are finite cooling preparations. For each,
\[
\pi_i(\beta)=\frac{q_i^\beta}{\sum_jq_j^\beta}.
\]
The certified slow eigenfunction, normalized as in Part II, gives
\[
-0.0937454280<b_h<-0.0937454279,
\qquad
-0.0938488357<b_c<-0.0938488356.
\]
Therefore
\[
|b_c|-|b_h|>0.0938488356-0.0937454280=0.0001034076>0.
\]
A stricter enclosure is recorded in the certificate. The gap is small but
strict. A small counterexample still refutes a universal implication.

\paragraph{Two witnesses, two sets of data.}
The focused paper uses a companion with rational slow eigenvalue and
rational eigenfunction. It changes two rates. Its formal proof and its
numerical values must not be silently substituted for the primary witness
above. The shared mechanism makes them related examples, not identical data.

\clearpage
\section*{Lesson 19. Work through a triplet test instead of treating it as a black box}
\addcontentsline{toc}{section}{19. Work through a triplet test}
A four-state system has four three-state subsets: $123$, $124$, $134$, and
$234$. For an ordered triplet $i<j<k$, define
\[
A=a_{ik},\quad B=a_{jk},\quad C=a_{ij},\quad
 t=q_j/q_i,\quad u=q_k/q_j.
\]
The letters are local abbreviations for this triplet. They do not introduce
new rates or fit effective ones. The original-rate predicate in Part II is
\begin{align*}
\kappa&=\frac{A+B}{C(1+t)},&G&=At-B,\\
F&=A(1+tu/2)-B(1+u/2)-C(1-t),\\
P&=(\kappa\geq1\text{ and }F>0)
\text{ or }(\kappa\leq1\text{ and }G>0).
\end{align*}
A predicate is simply a statement that is either true or false for the
specified data. The strict and non-strict inequality signs are part of the
definition. At $\kappa=1$, both branches must be considered.

\paragraph{Compute triplet $134$ by hand.}
Here $A=249/10$, $B=251/25$, $C=28$, $t=2/5$, and $u=1/160$.
First,
\[
G=\frac{249}{10}\frac25-\frac{251}{25}
=\frac{249-251}{25}=-\frac2{25}<0.
\]
Next,
\[
\kappa=\frac{249/10+251/25}{28(1+2/5)}
=\frac{1747/50}{196/5}=\frac{1747}{1960}<1.
\]
The applicable lower branch therefore fails because $G$ is negative.
For completeness, use $tu=1/400$ to calculate
\begin{align*}
F&=\frac{249}{10}\frac{801}{800}
-\frac{251}{25}\frac{321}{320}-28\frac35\\
&=\frac{199449-80571-134400}{8000}
=-\frac{7761}{4000}<0.
\end{align*}
The other three triplets are checked in the same manner in the exact table.
One failed triplet is not enough; the premise requires all four to fail.

\paragraph{Why isolated triplets are not principal submatrices.}
If a state is removed, its outgoing transitions from the retained states
must also be removed. The diagonal entries must be recomputed so the
three-state columns sum to zero. Merely deleting one row and column leaves
probability leaking to a state no longer in the model. This distinction is
separate from the predicate's use of original off-diagonal rates.

\clearpage
\section*{Lesson 20. Simplify the branch logic, including the equality case}
\addcontentsline{toc}{section}{20. Simplify the branch logic}
Complicated-looking logical tests can sometimes be converted into a small
set of inequalities. Define
\[
H=A+B-C(1+t),\qquad d=\frac{1-t}{1+t},
\qquad h=\frac2{1+t}+\frac u2.
\]
Because energies are distinct and ordered, $0<t<1$. Thus $d>0$ and $h>0$.
Also $C(1+t)>0$, so the sign of $\kappa-1$ is the sign of $H$.
A direct expansion gives
\[
F=hG+dH.
\]
To verify it, the coefficients of $A$, $B$, and $C$ on the right are,
respectively,
\[
\frac{2t+1-t}{1+t}+\frac{tu}{2}=1+\frac{tu}{2},
\quad\frac{-2+1-t}{1+t}-\frac u2=-1-\frac u2,
\quad-(1-t).
\]
These are exactly the coefficients in $F$.

\paragraph{Now handle the three cases.}
If $H>0$, only the upper branch can establish $P$. Its failure means
$F\leq0$. Since $F=hG+dH$ with $dH>0$, this forces $G<0$ as well.
If $H<0$, failure means $G\leq0$, and then $F=hG+dH<0$.
If $H=0$, both branches are eligible, but $F=hG$ with $h>0$; the strict
sign tests agree. Their joint failure means $F\leq0$ and $G\leq0$.
Conversely, if both are nonpositive, neither strict branch test can succeed.
Therefore, including equality cases,
\[
\neg P\quad\Longleftrightarrow\quad F\leq0\text{ and }G\leq0.
\]
The symbol $\neg$ means ``not.''

\paragraph{The geometric payoff.}
For fixed weights, $F$ and $G$ are linear functions of the rates. A condition
like $F\leq0$ describes a half-space. Intersecting the eight half-spaces from
four triplets gives a polyhedral cone: a set closed under addition and
multiplication by nonnegative numbers. Intersect it with strictly positive
rates to obtain the physical domain considered here. The word ``cone'' does
not assume a three-dimensional picture; the rate space has six coordinates.

This reformulation helps exact searches and structural proofs. It does not
replace the original predicate with a convenient approximation. The equality
above proves equivalence. Boundary cases remain part of the mathematical
problem even when the displayed witness lies strictly inside the domain.

\clearpage
\section*{Lesson 21. Certify a spectral calculation rather than trust its decimal output}
\addcontentsline{toc}{section}{21. Certify a spectral calculation}
A numerical eigensolver is excellent for finding candidate eigenvalues and
vectors. A certificate must additionally establish what those approximations
represent. Is the selected eigenvalue truly the slow one? Is it simple? Is
the overlap gap larger than all numerical uncertainty? The primary witness
answers these questions with exact algebra and rational enclosures.

\paragraph{The characteristic polynomial.}
The determinant $\det(M-xI)$ is a polynomial in $x$. Its zeros are the
eigenvalues. A determinant tests whether the columns are linearly dependent;
if $\det(M-xI)=0$, there is a nonzero vector $v$ with $(M-xI)v=0$.
Because $M$ is rational, its characteristic polynomial has rational
coefficients. The stationary eigenvalue supplies a factor $x$, leaving a
cubic. After clearing denominators, the primary witness's cubic is
\[
\begin{split}
1273600000000x^3+530473008000000x^2\qquad\qquad\quad\\
{}+35732953114072400x+510539126345027563.
\end{split}
\]
Its three roots lie separately in $(-20,-19)$, $(-60,-59)$, and
$(-337,-336)$. Exact sign checks at the endpoints establish a root in each
interval. A cubic has only three roots counting multiplicity, so these
three distinct intervals account for all roots and establish simplicity.
The root in $(-20,-19)$ is the slow nonzero one.

\paragraph{An eigenvector from cofactors.}
For any square matrix $A$, its adjugate satisfies
$A\operatorname{adj}(A)=\det(A)I$. At a simple eigenvalue $x$, a nonzero
column of $\operatorname{adj}(M^\top-xI)$ therefore lies in the nullspace
of $M^\top-xI$. Its entries are polynomials in $x$. The certificate verifies
that the chosen column is nonzero and normalizes it without dividing by an
interval containing zero. This avoids treating a rounded eigenvector as exact.

\paragraph{Thermal roots can be enclosed by integer powers.}
Suppose $r=q^{1/m}>0$, where $q$ is positive rational and $m$ is a positive
integer. If positive rational numbers $l<u$ satisfy $l^m<q<u^m$, then
$l<r<u$ because positive integer powers are strictly increasing. All
comparisons can be reduced to integer arithmetic by clearing denominators.
This is how the small preparation exponents can be handled rigorously.

\paragraph{Finally propagate the enclosures.}
Enclose the eigenvalue, evaluate the cofactor polynomials with interval
arithmetic, enclose each Gibbs weight, normalize by a strictly positive
partition-function interval, and sum the products. The resulting overlap
intervals must be disjoint in the required absolute-value order. If they
are not, increase precision or use a sharper method. A printed decimal
gap by itself is not the certificate; the outward bounds and their checked
construction are.

\clearpage
\section*{Lesson 22. Why three-state checks can miss four-state dynamics}
\addcontentsline{toc}{section}{22. Why local triplet checks can miss a collective mode}
A triplet test examines a restricted collection of rates. The slow
eigenfunction of the full four-state generator solves a simultaneous system
involving every state. Removing one state changes both the equations and
the probability-conserving diagonals. There is no general reason for the
full-system eigenfunction to be a collection of independent triplet
eigenfunctions pasted together.

\paragraph{A linear-algebra analogy.}
Knowing several smaller blocks of a matrix does not allow one to assign the
whole matrix's eigenvectors independently within those blocks. The
cross-block entries couple the coordinates. Similarly, a four-state mode
can distribute positive and negative weights among all four states so that
thermal reweighting produces a collective cancellation absent from any
isolated triplet. This explains why a local-to-global implication needs a
proof; it is not automatically supplied by the small subsystems.

\paragraph{What the counterexample does and does not say.}
It demonstrates failure of the stated universal finite-temperature
inference. It does not invalidate a distinct theorem with a zero-temperature
bath or stronger hypotheses. It also does not say that triplet information
is useless. The rest of the report derives substantial restrictions that
survive all those failed triplet tests. Ordinary cooling can survive while
strong effects and certain heating reversals remain excluded.

\paragraph{Temperature changes a weighted cancellation.}
The slow overlap is
\[
b(\beta)=\frac{w_1+w_2e^{-\beta E_2}
+w_3e^{-\beta E_3}+w_4e^{-\beta E_4}}{Z(\beta)}
\]
when $E_1=0$. The entries of $w$ need not all have the same sign, because
its bath mean is zero. Varying $\beta$ changes both the signed numerator
and its normalization. An intuitive story that only watches the highest
energy population misses the redistribution among the other three states.
The derivative is governed by the full covariance from Lesson 9.

\paragraph{The controlled-family idea.}
A single exact witness proves existence. A parameterized family is more
informative about mechanism: vary one parameter while retaining admissible
rates and triplet failures, then locate where the covariance criterion
changes sign. Part II supplies such a family with an exact onset threshold.
The family is a way to explain how the effect turns on, not a license to
assume that every nearby matrix has the same properties without checking.

\paragraph{Checkpoint.}
The logical conclusion is narrow and strong: there exists a valid
four-state example satisfying the full premise and violating the claimed
conclusion. Universal claims fail that way. Claims about prevalence,
experimental ease, or maximum possible benefit require additional work.

\clearpage
\section*{Lesson 23. Prefix contrasts, cones, and ordered dynamics}
\addcontentsline{toc}{section}{23. Prefix contrasts, cones, and ordered dynamics}
To understand what the triplet restrictions still control, it helps to
change basis in the space of observables. In this lesson let $q=\pi_b$ be
normalized; only ratios matter in the formulas. Define $Q_k=q_1+\cdots+q_k$
and a prefix contrast
\[
h^{(k)}_i=\begin{cases}
1,&i\leq k,\\
-Q_k/q_{k+1},&i=k+1,\\
0,&i>k+1.
\end{cases}
\]
It contrasts the first $k$ states against the next state. Its weighted
bath mean is $Q_k-q_{k+1}(Q_k/q_{k+1})=0$.
For example, with four states,
\[
h^{(1)}=(1,-q_1/q_2,0,0)^\top,\quad
h^{(2)}=(1,1,-Q_2/q_3,0)^\top.
\]
These are signed observables, not probability vectors.

\paragraph{Why these particular vectors are useful.}
If $j<k$, the vector $h^{(k)}$ equals one wherever $h^{(j)}$ is nonzero.
Thus their weighted inner product is the bath mean of $h^{(j)}$, which is
zero. They form an orthogonal basis of the bath-mean-zero observables.
In this basis, certain rate inequalities become entrywise sign conditions
on a smaller matrix. A matrix with nonnegative off-diagonal entries is
called a Metzler matrix. Adding a sufficiently large scalar multiple of
the identity turns it into a nonnegative matrix without changing its
eigenvectors. The Perron argument then supplies a nonnegative eigenvector
for its largest eigenvalue. Part II checks the needed hypotheses and
boundary cases instead of assuming them from the name of the theorem.

\paragraph{From the triplet inequality to conductance order.}
Define reduced conductances $k_{ij}=M_{ij}/q_i=M_{ji}/q_j$. Then
\[
G_{ijk}=a_{ik}\frac{q_j}{q_i}-a_{jk}
=q_j(k_{ik}-k_{jk}).
\]
Because $q_j>0$, the condition $G_{ijk}\leq0$ says exactly
$k_{ik}\leq k_{jk}$. The contrast representation exploits this ordering.
For a simple slow mode it allows an orientation
$w=\sum_k x_k h^{(k)}$ with $x_k\geq0$, not all zero.

\paragraph{Why this excludes a strong cooling zero.}
Consider one prefix contrast at a hotter preparation $\beta<\beta_b$.
For each $i\leq k$,
\[
\frac{\pi_i(\beta)}{\pi_{k+1}(\beta)}
=e^{\beta(E_{k+1}-E_i)}
<e^{\beta_b(E_{k+1}-E_i)}=\frac{q_i}{q_{k+1}}.
\]
Summing gives $\langle h^{(k)}\rangle_{\pi(\beta)}<0$.
Every nonnegative combination above therefore has strictly negative
overlap on the cooling side. It cannot cross zero there. Yet its negative
value may first move farther from zero and later return toward the bath:
a fixed sign does not imply monotonic magnitude. This is the space left
for ordinary cooling reversals.

\paragraph{A distinct order concept.}
Stochastic monotonicity asks whether dynamics preserve order of increasing
observables or cumulative probabilities. It is not the same statement as
nonnegativity of prefix-contrast coordinates. Part II derives the
appropriate adjacent-difference criterion for that stronger order property
and states exactly which observables and divergences it controls.

\clearpage
\section*{Lesson 24. Derive the sharp one-quarter amplitude bound}
\addcontentsline{toc}{section}{24. Derive the sharp one-quarter amplitude bound}
Under the full four-state triplet premise and a simple slow mode, Part II
proves the structural inequalities
\[
w_1\geq w_2,\qquad w_1+w_2\geq2w_3,\qquad w_4\leq0
\]
in the orientation of Lesson 23. The difficult rate-to-eigenfunction
argument uses a reset subtraction, a nonnegative contrast matrix, and a
Schur complement. We now derive the amplitude bound from these inequalities
step by step. This makes clear which part is structural input and which
part is elementary thermal averaging.

\paragraph{First condition on the lower three states.}
Let $m(\beta)$ be the Gibbs average of $w$ conditional on states $1,2,3$.
Write their three entries as a constant plus
\[
u(1,-1,0)+v(1,1,-2),\qquad
u=(w_1-w_2)/2\geq0,\quad
v=(w_1+w_2-2w_3)/6\geq0.
\]
This decomposition is checked by subtracting the first two coordinates
and by adding the first two minus twice the third. Those operations give
$2u$ and $6v$, respectively.

Set $E_1=0$, $a=E_2$, $b=E_3$, $t=e^{-\beta a}$ and $y=e^{-\beta b}$.
The conditional probabilities are $(p_1,p_2,p_3)=(1,t,y)/(1+t+y)$.
Quotient differentiation gives
\[
(p_1-p_2)'=\frac{2at+by-(b-a)ty}{(1+t+y)^2}>0.
\]
Indeed $b-(b-a)t\geq a>0$ for $0<t\leq1$. Also
$(p_1+p_2-2p_3)'=-3p_3'>0$, since state $3$ has the highest energy
in the conditional distribution. Nonnegative coefficients $u,v$ now imply
$m'(\beta)\geq0$.

\paragraph{Restore the fourth state.}
Let $P(\beta)=\pi_1+\pi_2+\pi_3$. It increases with $\beta$, and for
finite positive $\beta$ it lies strictly between $3/4$ and $1$.
The lower bound follows because each of the first three Gibbs weights
exceeds the fourth weight. The full overlap is
\[
b(\beta)=P(\beta)m(\beta)+(1-P(\beta))w_4.
\]
For $0<\beta_h<\beta_c<\beta_b$, put $r=P_c/P_h$. Then $1<r<4/3$.
Subtracting $rb_h$ from $b_c$ cancels the common lower-state weight:
\begin{align*}
b_c-rb_h
&=P_c(m_c-m_h)+(1-r)w_4\\
&\geq0.
\end{align*}
Both terms are nonnegative: $m_c\geq m_h$, while both $1-r$ and $w_4$
are nonpositive. Lesson 23 gives $b_h,b_c<0$. Replacing them by minus their
absolute values reverses the intuitive direction, yielding
\[
|b_c|\leq r|b_h|<\frac43|b_h|.
\]
Divide by the positive $|b_c|$ and rearrange:
\[
\frac{|b_h|}{|b_c|}>\frac34,
\qquad A=1-\frac{|b_h|}{|b_c|}<\frac14.
\]
That is the one-quarter bound. It applies to the stated amplitude measure
under the stated full premise, not to every four-state thermal model.

\paragraph{Why ``sharp supremum'' does not mean an attained maximum.}
A supremum is the least upper bound. The numbers $1/4-1/m$ approach $1/4$
without equaling it at any finite $m$. In the same way, the constructed
admissible families approach $25\%$ while every finite member stays below.
A limiting expression with a zero weight or infinite energy need not itself
be admissible; the proof must verify a sequence of finite admissible members.
The report supplies that verification, and distinguishes its uniformly
gapped, strict-inequality family from the different family used in Lean.
The corresponding asymptotic KL reduction approaches
$1-(3/4)^2=7/16$, not $1/4$.


\paragraph{One level deeper: what the Schur complement actually does.}
The structural proof mentions eliminating a coordinate. Here is the algebra
behind that phrase. Partition a three-by-three eigenvalue equation into
two coordinates $y=(x_1,x_2)^\top$ and a third coordinate $z=x_3$:
\[
\begin{pmatrix}A&c\\d^\top&e\end{pmatrix}
\begin{pmatrix}y\\z\end{pmatrix}
=\lambda\begin{pmatrix}y\\z\end{pmatrix}.
\]
Here $A$ is a two-by-two matrix, $c,d$ are two-component columns, and $e$
is a scalar. The bottom row says $d^\top y+ez=\lambda z$.
If $\delta=\lambda-e>0$, solve it as $z=d^\top y/\delta$.
Substituting into the top rows gives
\[
\left(A+\frac{cd^\top}{\delta}\right)y=\lambda y.
\]
Thus the effective two-by-two matrix includes the indirect coupling through
the eliminated coordinate. The operation is exact; it is not dropping the
third coordinate or assuming that it is small. If $c,d$ have nonnegative
entries and $\delta>0$, the additional coupling has nonnegative entries too.

Why not immediately divide? Because $\delta>0$ must first be established.
For the symmetric version of a matrix, the largest eigenvalue is at least
any diagonal entry: evaluate its quadratic form on the corresponding unit
coordinate vector and compare with the largest eigenvalue. Equality is a
special case. In the report's proof, equality would make that coordinate
vector a top eigenvector. Simplicity would then contradict a required
nonzero first coordinate. This supplies strictness, permitting division.
The zero-first-coordinate case is treated separately before this step.

Finally, for an effective matrix with entries $b_{ij}$, let
$r=x_2/x_1$ when $x_1>0$. Its two eigenvalue equations are
$b_{11}+b_{12}r=\lambda$ and $b_{21}+b_{22}r=\lambda r$.
Eliminate $\lambda$ to obtain
\[
b_{21}+r(b_{22}-b_{11})-b_{12}r^2=0.
\]
The structural proof evaluates this polynomial at a proposed lower bound
for $r$. Rate inequalities control its sign, and the positive-root
structure controls which side of the root the proposed bound lies on.
This is how a statement about rates becomes a statement about the shape
of the slow eigenfunction. The report also treats vanishing off-diagonal
entries, where a positive-root shortcut alone would be insufficient.

\clearpage
\section*{Lesson 25. From plotting a curve to classifying every preparation}
\addcontentsline{toc}{section}{25. From plotting a curve to classifying every preparation}
A plot of $b(\beta)$ suggests where extrema occur, but a finite collection
of samples cannot rule out an extra turn between sample points. A global
classification needs a bound on how complicated the curve can be and a
finite list of decisive tests.

\paragraph{Remove a redundant coordinate.}
For four distinct energies with $E_1=0$, the ground-state probability has
$\pi_1'>0$ at every finite $\beta\geq0$. Write
\[
g=\pi_1',\qquad u=\frac{\pi_2'}g,\qquad
v=\frac{\pi_3'}g.
\]
Since the derivative probabilities sum to zero,
$\pi_4'=-g(1+u+v)$. Consequently
\[
b'=g\{A+Bu+Cv\},\qquad
A=w_1-w_4,\ B=w_2-w_4,\ C=w_3-w_4.
\]
Because $g>0$, the sign of $b'$ is exactly the sign of
$\Phi=A+Bu+Cv$. Instead of looking at four derivatives separately, we
study a planar curve $(u,v)$ and a line $A+Bu+Cv=0$.

\paragraph{What the determinant calculation establishes.}
A Wronskian is a determinant whose columns are functions and whose rows
are successive derivatives. For two functions it is
$W(f,g)=fg'-f'g$. For more functions the determinant records higher-order
independence. In this problem the exact exponential structure of the
Gibbs probabilities makes relevant Wronskians factor into quantities of
known sign. Part II derives the factorization. Its consequences here are
$u'<0$ and strict convexity of $v$ as a function of $u$.

Strict convexity means that the curve bends consistently in one direction.
A line can meet such a curve at at most two points in this setting, with
multiplicity counted appropriately. More specifically, if
$s=dv/du$, then
\[
\Phi'=u'(B+Cs).
\]
The report gives an explicit monotone formula for $s$. Thus $\Phi$ has at
most one interior extremum unless it is a special constant case. To decide
whether $b'$ takes both signs on a bath side, it is enough to test $\Phi$
at that side's endpoints, its endpoint limits, and the possible single
interior extremum. Equality cases are retained, rather than assigned a
sign from a plotting tolerance.

\paragraph{Tangency versus reversal.}
If $\Phi$ merely touches zero and returns to the same sign, $b'$ does not
change sign. A stationary point of $b$ has occurred without a reversal of
monotonicity. If $b$ itself crosses zero away from the bath, that concerns
a strong preparation instead. These are two different root problems.
The classification separates them and uses monotone intervals to decide
where nonbath zeros of $b$ exist.

\paragraph{The advanced generalization.}
In higher dimension, exponential functions form a system with controlled
zero counts. The report derives a bound of $n-2$ on stationary points of a
nonconstant thermal overlap, with multiplicities treated in the relevant
interval, and constructs the matching number of strong preparations.
The four-state planar picture is a low-dimensional window into that
Chebyshev-system argument. The terminology comes after the concrete curve,
not in place of understanding it.

\clearpage
\section*{Lesson 26. Repeated slow eigenvalues and why one overlap is no longer enough}
\addcontentsline{toc}{section}{26. Repeated slow eigenvalues}
Suppose the slow eigenvalue $\lambda<0$ has two independent eigenvectors.
Every combination of them decays with the same exponential. Choosing one
of those vectors is a coordinate choice, not a preferred physical mode.
A preparation may have zero coefficient in one chosen direction but a
nonzero coefficient in the other. Calling that a strong effect would be
incorrect.

\paragraph{Use the entire projection.}
In symmetric coordinates from Lesson 15, let $P_\lambda$ be the orthogonal
projection onto the whole slow eigenspace. For a preparation, form
\[
y(\beta)=D^{-1/2}(\pi(\beta)-\pi_b),
\qquad A_\lambda(\beta)=\|P_\lambda y(\beta)\|_2^2.
\]
The subscript $2$ denotes ordinary Euclidean length. The leading Pearson
chi-square divergence is
$e^{2\lambda t}A_\lambda(\beta)$. Leading KL is one half of it. A strong
preparation must have $P_\lambda y=0$, not merely one zero coordinate.
For other norms, the shape of the projected vector matters, so a single
quadratic coefficient need not decide every distance comparison.

\paragraph{A reset process makes the issue visible.}
Let
\[
M=r(\pi_b\mathbf1^\top-I),\qquad r>0,
\]
where $\mathbf1$ is the all-ones column. Since $\mathbf1^\top p=1$,
$Mp=r(\pi_b-p)$, so
\[
p(t)=\pi_b+e^{-rt}(p(0)-\pi_b).
\]
Every zero-total deviation is an eigenvector with eigenvalue $-r$.
There are $n-1$ independent slow directions. All deviations shrink by the
same factor. There is no need to choose a distinguished vector inside
this eigenspace, and doing so cannot reveal a physical effect absent from
the full projection.

\paragraph{Where the harder classification begins.}
Part II studies repeated slow spaces and, under the full triplet premise,
classifies the four-state boundary as a reset plus mixing on one adjacent
pair. It then proves the appropriate all-time convex-divergence ordering.
This is a separate structural result, not something that follows merely
from the word ``degenerate.'' Degeneracy means repeated eigenvalues;
it does not mean singular probabilities or an invalid generator.

\paragraph{Checkpoint.}
For a vector with two coordinates $(3,4)$, one coordinate can be made zero
by rotating axes, while its length remains $5$. The same principle explains
why a basis-dependent zero overlap is inadequate in a repeated eigenspace.
The invariant object is the whole projected deviation.

\clearpage
\section*{Lesson 27. Numerical evidence, exact arithmetic, and interval proof}
\addcontentsline{toc}{section}{27. Numerical evidence, exact arithmetic, and interval proof}
Computers do several distinct jobs in this project. A search proposes
parameters. A numerical solver estimates their behavior. Exact algebra
checks identities. Interval arithmetic encloses uncertain real quantities.
A proof assistant checks a formal argument. These jobs reinforce one
another, but one does not automatically substitute for another.

\paragraph{Floating-point arithmetic.}
Ordinary numerical software stores finitely many binary digits. Most real
numbers, and many simple decimal fractions, cannot be represented exactly.
Each operation can introduce rounding. This is usually harmless for
exploration, but the sign of a tiny difference between two nearly equal
numbers can be sensitive. Subtracting two quantities around $1$ to obtain
a result around $10^{-12}$ demands a much more careful error analysis
than merely reporting twelve displayed digits.

\paragraph{Exact rational arithmetic.}
Integers and ratios of integers can be stored and manipulated exactly.
For instance,
$249/25-251/25=-2/25$ is an exact negative result. Exact arithmetic is ideal
for the rate construction, polynomial coefficients, and cleared inequalities.
It does not imply that every answer must be rational: an irrational
polynomial root can be specified exactly by its polynomial and a rational
isolating interval.

\paragraph{Interval arithmetic.}
Replace an unknown real number by a closed interval guaranteed to contain
it. For $x\in[a,b]$ and $y\in[c,d]$,
\[
x+y\in[a+c,b+d],\qquad x-y\in[a-d,b-c].
\]
A product is enclosed by the minimum and maximum of the four endpoint
products. Division requires that the denominator interval not contain zero.
With machine arithmetic, endpoints must be rounded outward so the true
answer cannot escape. Functions such as exponential and logarithm likewise
require validated enclosures, not ordinary rounded evaluations casually
wrapped in brackets.

\paragraph{Dependency can make a valid enclosure too wide.}
If $x\in[0,1]$, a naive interval evaluation of $x-x$ gives $[-1,1]$,
although the exact expression is zero. The enclosure is valid but weak
because the same variable appeared twice. Algebraic simplification,
smaller boxes, or sharper bounds can restore useful precision. A failed
interval inequality may mean only that the enclosure is too broad.

\paragraph{Why random perturbations do not certify a box.}
Testing ten thousand parameter points leaves infinitely many others
untested. To prove an inequality on an entire box, bound it on that box;
if the bound is inconclusive, split the box and repeat. A complete
branch-and-bound certificate must account for every subbox, including
ones near difficult boundaries. A pending queue, a timeout, or an enclosure
error must remain visible. Success on most of the domain is not a proof
on all of it.

\paragraph{Negative controls are part of credibility.}
A checker should reject intentionally corrupted data and deliberately
false claims. Otherwise a script that always reports success could look
impressive. The project includes such controls, alongside independent
formulations and source-bound receipts. They improve confidence in the
implementation without becoming external human peer review.

\clearpage
\section*{Lesson 28. From late-time asymptotics to useful finite-time questions}
\addcontentsline{toc}{section}{28. From asymptotics to finite-time questions}
A strict slow-amplitude reversal proves an eventual comparison. It does
not tell us whether the crossing occurs before the deviations become
experimentally negligible. This motivates a separate question: can a
specified model show a substantial advantage over a stated time interval
and at stated, nonzero distance thresholds?

\paragraph{Specify the observable and the clock.}
Let $D_x(t)=D_{\mathrm{KL}}(p_x(t)\Vert\pi_b)$ for preparation $x$.
A distance advantage of at least $10\%$ means
$D_h(t)\leq0.9D_c(t)$ on the specified interval. A floor
$D_c(t)\geq10^{-3}$ prevents the claim from being supported solely by
arbitrarily tiny distances. Neither statement is a target-time claim.
Define a threshold time instead by
\[
t_x(\delta)=\inf\{t\geq0:D_x(t)\leq\delta\}.
\]
A time saving of at least $5\%$ means $t_h(\delta)\leq0.95t_c(\delta)$.
Here $\delta$ is a distance threshold, not a temperature.

\paragraph{Monotonicity makes threshold roots usable.}
For the irreducible reversible dynamics studied in the physical section,
KL decreases strictly along a nonstationary trajectory. A threshold between
its initial value and zero has a unique crossing time. If a candidate time
$s$ satisfies $D_h(s)\leq\delta$, then $t_h(\delta)\leq s$.
Thus, taking $s=0.95t_c(\delta)$, a sufficient condition for the desired
time saving is
\[
D_h(0.95t_c(\delta))\leq D_c(t_c(\delta))=\delta.
\]
The equation explains the logic; the certificate must still handle the
unknown cold root and every allowed parameter vector.

\paragraph{Normalize time carefully.}
The physical study writes $u=t/\tau$, where
$\tau=-1/\lambda_2$ is the relaxation time of the chosen generator.
For a fixed parameter vector, both preparations share that same $\tau$,
so a fractional target-time comparison is unchanged by normalization.
Across a parameter box, however, $\tau$ itself changes. A dimensionless
window does not automatically specify one shared number of seconds for
all models. An attempt frequency is needed to turn the clock into seconds.

\paragraph{The certified physical-model benchmarks.}
The retained study establishes $D_h\leq0.9D_c$ and $D_c\geq10^{-3}$ for
$0.4\leq u\leq0.8$, and $t_h(\delta)\leq0.95t_c(\delta)$ for every
$\delta\in[10^{-3},10^{-2}]$, throughout its specified simultaneous
$\pm5\%$ parameter box. The perturbed coordinates are three energy gaps,
three shared barrier excesses, and two temperature excesses $T_x/T_b-1$.
The last phrase matters: perturbing a temperature excess is not identical
to perturbing the full temperature ratio. The eight continuous obligations
and their reductions are recorded in Part II.

\clearpage
\section*{Lesson 29. What the physical model adds, and what remains idealized}
\addcontentsline{toc}{section}{29. What the physical model adds}
The exact counterexample was designed to settle a mathematical inference.
The physical-model study asks a different question using a constrained
thermal network. Its rate law has Arrhenius form
\[
a_{ij}=\nu_0\exp[-\beta_b(B_{ij}-E_j)],\qquad B_{ij}=B_{ji}.
\]
Here $\nu_0$ is an attempt frequency and $B_{ij}$ is a shared saddle energy.
The exponent uses the departure energy $E_j$, consistent with the
 destination-first index convention. A higher barrier above the departure
state suppresses that transition.

\paragraph{Detailed balance follows from the symmetric saddle.}
Multiply the rate by the unnormalized bath weight $q_j=e^{-\beta_b E_j}$:
\[
a_{ij}q_j
=\nu_0e^{-\beta_b(B_{ij}-E_j)}e^{-\beta_b E_j}
=\nu_0e^{-\beta_b B_{ij}}
=a_{ji}q_i.
\]
The departure energy cancels. This is a transparent physical construction
of reversible rates, rather than a separate balance constraint fitted
after the fact.

\paragraph{Shared parameters reduce arbitrary freedom.}
The study uses four ordered energies with three gaps and two interleaved
basins, $\{1,4\}$ and $\{2,3\}$. The saddle is
$B_{ij}=\max(E_i,E_j)+h_{ij}$, with one shared excess for each intrabasin
edge type and one for all interbasin edges. The six edges therefore do not
have six independently fitted barrier excesses. This geometry makes the
model more constrained and supplies a kinetic interpretation of the
collective relaxation.

\paragraph{Do not merge the two research claims.}
Two triplets in this physical example satisfy the original triplet test.
It is therefore not another witness to the all-triplet-failure conjecture.
The abstract counterexample and the robust physical example answer
different questions. Keeping them separate prevents a compelling finite-time
plot from being used as evidence for a premise it does not satisfy.

\paragraph{A realization is not an experiment.}
Part II gives an energy-conserving stochastic collision realization of the
rates and examines an explicit activated-intermediate construction. These
make assumptions more inspectable. They do not identify a particular
material, prove that a laboratory contact process obeys the idealization,
or include all costs of preparing a hotter initial state. A demonstrated
reduction in the model's relaxation time is not automatically a net energy
or engineering advantage for a device. That latter question requires an
account of preparation, coupling, measurement, and operating constraints.

\paragraph{A productive next layer of work.}
To connect the model to an experiment, specify what each state represents,
how rates are measured, how the bath approximation is maintained, and which
observable supplies the stopping criterion. Then compare model predictions
with measured uncertainty. The mathematical certificate is valuable because
it makes one link in that chain explicit; it does not claim the other links
have already been established.

\clearpage
\section*{Lesson 30. What a formal proof checks, and how to read its scope}
\addcontentsline{toc}{section}{30. Formal proof and reproducibility}
A proof assistant checks a mathematical argument written in a formal
language against explicit definitions and inference rules. The central
question becomes: what exact proposition was checked? A successful build
of one theorem does not imply that every sentence in a surrounding report
has been formalized.

\paragraph{Start with the definitions.}
For this project, a formal counterexample must define the actual generator,
its bath weights, the preparation distributions, the triplet predicate,
and the chosen distance. If any of these definitions differs from the
informal claim, a formally correct theorem could still answer the wrong
question. Source-to-statement correspondence is therefore a separate
review obligation. The same applies to exact computational certificates.

\paragraph{The covered results.}
The Lean companion verifies the rational-mode counterexample, an
initial-to-eventual Pearson chi-square reversal under its master equation,
and the sharp four-state amplitude supremum. It uses a particular sharpness
family. The report's stronger all-strict, uniformly gapped family has a
separate analytic and computational justification; it is not made formal
merely because a related family is formalized. The physical interval
certificates are also a distinct verification route.

\paragraph{What reproducibility should expose.}
A useful reproduction record identifies the inputs, source version, command,
output, and interpretation. A cryptographic hash is a compact identifier
of file contents. Matching hashes tie an old result to the same bytes;
they do not prove those bytes implement the intended mathematics. A changed
source invalidates that simple correspondence until the change has been
reviewed. This is why the accompanying guide tracks source-bound proof
reviews separately from passing replay logs.

\paragraph{A practical reading order for a theorem.}
First read its hypotheses, including positivity, simplicity, temperature
range, and chosen distance. Then read the conclusion literally: existence,
all-time ordering, eventual ordering, or a quantitative bound. Next identify
which definitions give those words mathematical content. Finally inspect
the proof and its dependencies. Running the command comes after knowing
what a successful command is meant to establish.

\paragraph{How much confidence is justified?}
Independent exact arithmetic, alternate spectral arguments, negative
controls, interval bounds, and formal checking each address different
failure modes. Their agreement is useful evidence. None should be
represented as external peer review or experimental confirmation. The
report preserves those distinctions so that another researcher can decide
what to reuse and what to check independently.

\clearpage
\section*{Lesson 31. Exercises with complete worked answers}
\addcontentsline{toc}{section}{31. Exercises with complete worked answers}
These problems revisit the dependency chain. Try each calculation before
reading its answer. Their purpose is to reveal whether a symbol has acquired
a concrete meaning, not to test speed.

\paragraph{Exercise 1: normalize and compare.}
Given weights $(1,1/2,1/8)$, find the probabilities and the ratio of the
third to the first. \emph{Answer:} the sum is $13/8$, so the probabilities
are $(8/13,4/13,1/13)$. Their ratio is $(1/13)/(8/13)=1/8$, equal to the
ratio of the original weights. Normalization changes entries but cancels
from pair ratios.

\paragraph{Exercise 2: translate a temperature ordering.}
If $T_h=4T_b$ and $T_c=2T_b$, find the inverse temperatures in bath units.
\emph{Answer:} $\beta_h=1/4$, $\beta_c=1/2$, and $\beta_b=1$.
Both preparations lie on the cooling side, with $0<1/4<1/2<1$.
The preparation called ``cold'' is still hotter than the bath.

\paragraph{Exercise 3: construct both directions.}
Let $q=(1,1/4)$ and choose downward rate $a_{12}=2$. Find the upward rate,
generator, equilibrium, and nonzero eigenvalue. \emph{Answer:}
$a_{21}=2(1/4)=1/2$, so
\[
M=\begin{pmatrix}-1/2&2\\1/2&-2\end{pmatrix},
\qquad \pi_b=(4/5,1/5)^\top.
\]
The columns sum to zero, and the equilibrium flows are
$2(1/5)=(1/2)(4/5)=2/5$. One eigenvalue is zero; their sum is the trace
$-1/2-2=-5/2$, so the other is $-5/2$. Direct multiplication by $(1,-1)^\top$
checks the nonzero mode independently.

\paragraph{Exercise 4: verify a covariance sign.}
Use $p=(1/2,1/2)$, $E=(0,2)$, and $w=(1,-1)$.
\emph{Answer:} $\langle E\rangle=1$, $\langle w\rangle=0$, and
$\langle wE\rangle=-1$, so $\Cov(w,E)=-1$. For a fixed observable in a
Gibbs family passing through these probabilities, $b'=1$ there. The
pairwise form gives $(1/2)(1/2)(-2)(2)=-1$ as a second check.

\paragraph{Exercise 5: distinguish sign and magnitude.}
Suppose certified intervals give $b_h\in[-0.21,-0.20]$ and
$b_c\in[-0.24,-0.23]$. Prove a strict magnitude advantage and bound its
fraction conservatively. \emph{Answer:} $|b_h|\leq0.21<0.23\leq|b_c|$.
Thus $1-|b_h|/|b_c|\geq1-0.21/0.23=2/23$, about $8.7\%$.
Using interval midpoints would give an estimate, not this guaranteed bound.

\paragraph{Exercise 6: calculate the two-exponential crossing.}
For $f_h=e^{-t}+9e^{-10t}$ and $f_c=2e^{-t}+e^{-10t}$, determine their
ordering. \emph{Answer:} $f_h-f_c=e^{-10t}(8-e^{9t})$.
The prefactor is positive, so the difference is positive before
$t=\log8/9$, zero there, and negative after it. This is a complete sign
argument, with no plotting or root-finding tolerance.

\paragraph{Exercise 7: separate amplitude, divergence, and time.}
If a one-mode amplitude ratio is $0.8$, what is the late-time KL ratio?
\emph{Answer:} $0.8^2=0.64$, giving a $36\%$ asymptotic KL reduction.
With slow decay rate $r$, the leading target-time advance is
$\log(1/0.8)/r$. Its percentage depends on the reference target time.
No universal percentage time saving follows from the amplitude ratio alone.

\paragraph{Exercise 8: find the missing logical check.}
A program finds a four-state reversal, but checks only three of the four
triplets. Does it certify the counterexample? \emph{Answer:} no. The
antecedent requires all four predicates to fail. A single omitted triplet
could satisfy the test and invalidate the proposed contradiction even if
the reversal itself is real.

\paragraph{Exercise 9: understand a strict supremum.}
Every allowed model has $A<1/4$, while a sequence of allowed models has
$A_m\to1/4$. Can one claim an allowed model attains $25\%$?
\emph{Answer:} no. The two facts establish the supremum exactly, but the
strict bound rules out attainment. A limit outside the allowed domain
does not supply a missing maximizer.

\paragraph{Exercise 10: identify insufficient robustness evidence.}
A numerical study samples a parameter box, finds no violations, and reports
a $5\%$ guarantee everywhere. What is missing? \emph{Answer:} a proof
covering every point between samples and numerical error in the evaluated
quantities. A complete interval cover can provide such a proof if every
subbox is successfully certified and the implementation bounds are valid.
Random sample count alone cannot provide it.


\paragraph{Exercise 11: verify a prefix contrast.}
Use normalized weights $q=(4/10,3/10,2/10,1/10)$.
Find $h^{(2)}$ and its squared weighted length.
\emph{Answer:} $Q_2=7/10$ and $Q_2/q_3=7/2$, so
$h^{(2)}=(1,1,-7/2,0)$. Its weighted mean is
$4/10+3/10-(2/10)(7/2)=0$. Its squared length is
$4/10+3/10+(2/10)(49/4)=63/20$. A signed contrast can have length larger
than one; it is not a probability vector.

\paragraph{Exercise 12: test a Schur elimination by substitution.}
Suppose an eigenvalue equation has $e=-5$, $\lambda=-2$,
$c=(1,2)^\top$, and $d=(3,4)^\top$. Find the correction to $A$.
\emph{Answer:} $\delta=3$ and
\[
\frac{cd^\top}{\delta}
=\frac13\begin{pmatrix}3&4\\6&8\end{pmatrix}.
\]
The bottom equation gives $z=(3x_1+4x_2)/3$.
Substituting $cz$ into the top equation yields exactly this correction.
This example checks the elimination algebra; no claim is made that an
arbitrary choice of $A$ gives a Markov generator or the proposed eigenvalue.

\paragraph{Exercise 13: detect an illegal division.}
An interval algorithm encloses a denominator by $[-0.01,0.02]$ and then
reports the quotient using the reciprocals of the endpoints.
\emph{Answer:} the step is invalid because the interval includes zero.
The reciprocal is unbounded near zero. One must establish a narrower
nonzero enclosure, split and analyze the domain with the singularity
handled, or reformulate the comparison without that division. A small
interval is not enough; its position relative to zero matters.

\paragraph{Exercise 14: turn a distance bound into a time bound.}
Assume a strictly decreasing $D_h(t)$, and let $D_c(t_c)=\delta$.
If $D_h(0.95t_c)\leq\delta$, show the time saving.
\emph{Answer:} by time $0.95t_c$ the hot trajectory has already reached
the threshold. Its first crossing therefore satisfies
$t_h\leq0.95t_c$. Subtracting from $t_c$ gives
$t_c-t_h\geq0.05t_c$, a saving of at least $5\%$ relative to $t_c$.
A pointwise comparison $D_h(t_c)<D_c(t_c)$ alone would establish only
$t_h<t_c$, without the quantitative $5\%$ margin.

\clearpage
\section*{Lesson 32. Map the foundations into the advanced report}
\addcontentsline{toc}{section}{32. Map the foundations into the advanced report}
The next part preserves the technical research exposition. It uses the
language now developed, but it is intentionally denser. Use the following
routes to move between explanation and proof rather than trying to absorb
every theorem in one uninterrupted reading.

\begin{longtable}{>{\raggedright\arraybackslash}p{0.25\textwidth}>{\raggedright\arraybackslash}p{0.65\textwidth}}
\toprule Question & Route through the report\\\midrule
\endhead
What exactly was contradicted? & Read the source predicate and
Section~\ref{sec:claimscope}; match each hypothesis to its witness before
reading the eigenvalue calculation. Revisit Lessons 1, 19, and 20.\\[5pt]
Why does the overlap imply relaxation reversal? & Read
Section~\ref{sec:cooling} with Lessons 14, 16, and 17. Track the remainder
and the simple-slow-mode assumption.\\[5pt]
Can I check the original witness independently? & Read
Theorem~\ref{thm:counterexample}, the cofactor appendix, and
Section~\ref{sec:separatecheck}. Reconstruct the rates as in Lesson 18 and
understand the enclosures as in Lesson 21.\\[5pt]
What creates the effect? & Read Section~\ref{sec:mechanism}, then the
contrast argument in Section~\ref{sec:orderedconductances}. Revisit the
covariance identity and prefix contrasts.\\[5pt]
Why is the amplitude bound exactly one quarter? & Read
Section~\ref{sec:fulltripletbound} after Lesson 24. Separate the structural
lemma, the conditional-average bound, and the construction proving
sharpness. They are three different proof obligations.\\[5pt]
What changes without the triplet premise? & Read
Section~\ref{sec:universal} for the four-state classification and
Section~\ref{sec:degenerateextension} for repeated slow spaces.
Lessons 25--26 explain the geometry and the invariant projection.\\[5pt]
How do exact algorithms decide a case? & Read
Section~\ref{sec:exactclassifier} and the later rational-input classifier.
Watch assumptions on energies, exact root isolation, equality handling,
and failure reporting.\\[5pt]
What survives at finite time? & Read Section~\ref{sec:finitechi} for a
specified pair and the finite-time order sections for their precise
hypotheses. Eventual order alone is not their proof.\\[5pt]
What is physically demonstrated by the model? & Read
Sections~\ref{sec:physicalresults}, \ref{sec:physicalboxcertificate}, and
\ref{sec:thermalcollision}. Keep the two experiments of Lessons 18 and
29 separate.\\[5pt]
What has been formalized? & Read Section~\ref{sec:lean} with Lesson 30.
Match named theorems to the stated definitions and scope.\\
\bottomrule
\end{longtable}

\paragraph{A compact symbol dictionary.}
\begin{longtable}{>{\raggedright\arraybackslash}p{0.22\textwidth}>{\raggedright\arraybackslash}p{0.68\textwidth}}
\toprule Symbol & Meaning and common trap\\\midrule
\endhead
$i,j,k,n$ & State indices and number of states; an index is a label, not a time.\\
$E_i$ & Energy of state $i$; only energy differences affect Gibbs probabilities.\\
$T,\beta$ & Temperature and inverse temperature; their orderings are opposite.\\
$\beta_b,\pi_b$ & Fixed bath inverse temperature and its Gibbs distribution.\\
$q_i,Q$ & Positive bath weights and their sum; some advanced sections explicitly
use normalized $q=\pi_b$ instead. Read the local definition.\\
$\pi(\beta),Z(\beta)$ & Gibbs preparation and its partition function.\\
$p(t)$ & Evolving probability vector; generally not itself a Gibbs preparation.\\
$a_{ij},M_{ij}$ & Off-diagonal rate from $j$ to $i$; columns, not rows,
conserve probability in this convention.\\
$M,L$ & Probability generator and observable generator $L=M^\top$.\\
$D,S$ & Diagonal bath-probability matrix and symmetric transform
$S=D^{-1/2}MD^{1/2}$.\\
$\lambda_2,v_2,w_2$ & Slow nonzero eigenvalue, right eigenvector, and left
eigenfunction; simplicity must be checked.\\
$b(\beta)$ & Slow overlap $w_2^\top\pi(\beta)$, with a fixed normalization.\\
$\Cov,\Var$ & Covariance and variance under the indicated probability distribution.\\
$F,G,\kappa,P_{ijk}$ & Exact triplet expressions and predicate; $P_{ijk}$
is a logical statement, not a transition matrix.\\
$\tau,u$ & Relaxation time and normalized time $u=t/\tau$ in the physical study;
the letter $u$ has different explicitly defined local uses elsewhere.\\
$\delta,t_x(\delta)$ & A distance threshold and its crossing time.\\
$O(\cdot),o(\cdot)$ & Controlled order of a remainder and a relatively vanishing
remainder, respectively; neither is a blanket finite-error certificate.\\
\bottomrule
\end{longtable}

\paragraph{The habit to carry forward.}
For every advanced statement, ask what varies and what stays fixed; which
quantity is ordered; whether the comparison is finite-time or asymptotic;
where strictness enters; and how equality or degeneracy is handled.
Those questions are often more useful than memorizing a long formula.
The technical part that follows supplies the proofs, exact data, source
correspondence, implementation details, and explicitly stated open problems.

% PHYSICAL_TEACHING_BRIDGE

\clearpage
\part*{Part II. Technical research report and exact certificates}
\addcontentsline{toc}{part}{Part II. Technical research report and exact certificates}
\noindent The abstract below summarizes the research results in their technical
language. All original numbered sections, theorems, equations, and proof text
are retained after this teaching part; the foundational examples above use
unnumbered displays so that those references keep their original numbering.
\par\bigskip
